\section*{Capítulo \ref{ch:b1:enzymes} --- Enzimas y catálisis bioquímica}

\begin{solution}{exo:b1:enzymes:1}
Rebaja la \omterm{prop:b1:enzymes:whatitdoes}{energía de activación} (la altura del \omterm{prop:b1:enzymes:whatitdoes}{estado de transición}) y, con
ella, la velocidad en ambos sentidos; deja sin cambios $\Delta G$, la
constante de equilibrio y la posición del equilibrio. En el perfil, el pico
baja y los dos extremos no.
\end{solution}

\begin{solution}{exo:b1:enzymes:2}
$K_m$: concentración de \omterm{def:b1:enzymes:enzyme}{sustrato} a la mitad de $V_{\max}$ (\unit{mol/L}).
$V_{\max}$: velocidad con \omterm{def:b1:enzymes:enzyme}{sustrato} saturante (\unit{mol/s}, o \unit{\micro
mol/min}). $k_{\text{cat}}$: moléculas de \omterm{def:b1:enzymes:enzyme}{sustrato} convertidas por segundo y
por \omterm{def:b1:enzymes:enzyme}{enzima} en saturación (\unit{s^{-1}}). $k_{\text{cat}}/K_m$: la eficacia,
la constante de velocidad de segundo orden a baja concentración de \omterm{def:b1:enzymes:enzyme}{sustrato}
(\unit{L.mol^{-1}.s^{-1}}).
\end{solution}

\begin{solution}{exo:b1:enzymes:3}
Sin \omterm{def:b1:enzymes:inhibitors}{inhibidor}: ordenada en el origen 1, luego $V_{\max} = 1$; corte con el
eje $x$ en $-0.5$, luego $K_m = 2$. Competitivo: $V_{\max} = 1$, corte con
el eje $x$ en $-0.25$, $K_m = 4$. No competitivo: ordenada 2,
$V_{\max} = 0.5$; corte con el eje $x$ en $-0.5$, $K_m = 2$.
\end{solution}

\begin{solution}{exo:b1:enzymes:4}
Regulación \omterm{def:b1:proteins:allostery}{alostérica} (por retroalimentación), milisegundos; \omterm{prop:b1:enzymes:control}{modificación
covalente} (fosforilación), de segundos a minutos; activación proteolítica de
un zimógeno, una sola vez y a demanda; control de la cantidad por expresión
génica y degradación, de minutos a horas.
\end{solution}

\begin{solution}{exo:b1:enzymes:5}
$v_0 = 50[S]/(0.1 + [S])$: 8,3, 25, 45,5 y \qty{49.5}{\micro mol/min}.
$0.9 = [S]/(0.1 + [S])$ da $[S] = 0.9\,K_m/0.1 = \qty{0.9}{mmol/L}$.
\end{solution}

\begin{solution}{exo:b1:enzymes:6}
$e^{20\,000/2580} = e^{7.75} = 2300$. Para $10^{10}$: $\Delta E_a =
RT\ln 10^{10} = 2.58\times 23.0 = \qty{59}{kJ/mol}$.
\end{solution}

\begin{solution}{exo:b1:enzymes:7}
$V_{\max} = \qty{0.12}{mol/L}$ por minuto en \qty{1}{mL}:
\qty{1.2e-4}{mol/min} $= \qty{2.0e-6}{mol/s}$ de producto; \omterm{def:b1:enzymes:enzyme}{enzima},
\qty{2e-12}{mol}: $k_{\text{cat}} = 2\times 10^{-6}/2\times 10^{-12} =
\qty{1e6}{s^{-1}}$. Eficacia $10^6/0.012 = \qty{8e7}{L.mol^{-1}.s^{-1}}$,
dentro de un factor diez del límite de difusión: casi todo encuentro con una
molécula de \omterm{def:b1:enzymes:enzyme}{sustrato} es productivo.
\end{solution}

\begin{solution}{exo:b1:enzymes:8}
$1 + 2/K_i = 2$: $K_i = \qty{2}{mmol/L}$. Para el \qty{90}{\%}:
$[S] = 9K_m^{\text{ap}}$: \qty{9}{mmol/L} sin el \omterm{def:b1:enzymes:inhibitors}{inhibidor} y
\qty{18}{mmol/L} con él.
\end{solution}

\begin{solution}{exo:b1:enzymes:9}
Inhibir el primer paso comprometido detiene toda la vía de golpe y no
malgasta intermediarios, que de otro modo se acumularían; el producto final
es la señal que importa --- su abundancia es exactamente la información de
que la vía ya no hace falta ---, mientras que el nivel de un intermediario
no dice nada sobre la demanda.
\end{solution}

\begin{solution}{exo:b1:enzymes:10}
Una proteasa activa en la \omterm{def:b1:cell-unit-of-life:cell}{célula} que la fabrica digeriría esa \omterm{def:b1:cell-unit-of-life:cell}{célula};
fabricada como zimógeno inactivo, es inofensiva hasta que la
enteropeptidasa, una \omterm{def:b1:enzymes:enzyme}{enzima} del revestimiento intestinal, corta el
tripsinógeno a tripsina en el duodeno, donde se quiere la digestión; la
tripsina activa después los demás zimógenos (una cascada). Como una traza de
tripsina puede iniciar la cascada antes de tiempo, el páncreas empaqueta en
sus gránulos un \omterm{def:b1:enzymes:inhibitors}{inhibidor} específico de la tripsina como seguro; el fallo de
ese dispositivo es la pancreatitis.
\end{solution}

\begin{solution}{exo:b1:enzymes:11}
Sin él: $3^3/(3^3 + 3^3) = 0.50$. Con él: $27/(166 + 27) = 0.14$: la
velocidad cae al \qty{28}{\%} de su valor. En una \omterm{def:b1:enzymes:enzyme}{enzima} de
Michaelis--Menten con la $K_m$ elevada de 3 a 5,5: de $3/6 = 0.50$ a
$3/8.5 = 0.35$, es decir, el \qty{70}{\%}. La \omterm{def:b1:proteins:allostery}{cooperatividad} hace que ese
mismo desplazamiento de la curva sea más del doble de eficaz cerca del punto
de trabajo: un interruptor y no un regulador de intensidad.
\end{solution}

\begin{solution}{exo:b1:enzymes:12}
Un catalizador desnudo (una superficie de platino, un ácido) acelera muchas
reacciones sin discriminar; una \omterm{def:b1:enzymes:enzyme}{enzima} acelera una reacción sobre un
\omterm{def:b1:enzymes:enzyme}{sustrato}, y la especificidad de su \omterm{def:b1:enzymes:enzyme}{centro activo} es ya información sobre lo
que la \omterm{def:b1:cell-unit-of-life:cell}{célula} quiere que se haga. La regulación añade una segunda capa: los
sitios alostéricos, los sitios de fosforilación y los péptidos de los
zimógenos le dicen a la \omterm{def:b1:enzymes:enzyme}{enzima} cuándo y dónde actuar, según el estado de la
\omterm{def:b1:cell-unit-of-life:cell}{célula}. El coste es una \omterm{def:b1:proteins:peptide}{proteína} grande (cientos de residuos para un centro
de una docena), un recambio más lento que el de los mejores catalizadores
inorgánicos, y la maquinaria para fabricarla, modificarla y destruirla; el
beneficio es una química que funciona solo donde y cuando resulta útil, que
es lo que es un metabolismo.
\end{solution}

\begin{solution}{pb:b1:enzymes:1}
\textbf{1.} $1/[S]$: 2, 1, 0,5, 0,2, 0,1, 0,05. $1/v_0$: 5,0, 3,0, 2,0,
1,4, 1,2, 1,1.
\textbf{2.} Cada paso de $1/[S]$ cambia $1/v_0$ en proporción: pendiente
$(5.0 - 1.1)/(2 - 0.05) = 2.0$; ordenada en el origen $1.0$.
\textbf{3.} $V_{\max} = 1/1.0 = \qty{1.0}{\micro mol/min}$; $K_m =
\text{pendiente}\times V_{\max} = \qty{2.0}{mmol/L}$.
\textbf{4.} $1.0\times 5/(2 + 5) = 0.714$: lo medido.
\textbf{5.} $\qty{10e-9}{mol/L}\times\qty{1e-3}{L} = \qty{1e-11}{mol}$;
$V_{\max} = \qty{1e-6}{mol/min} = \qty{1.67e-8}{mol/s}$;
$k_{\text{cat}} = 1.67\times 10^{-8}/10^{-11} = \qty{1670}{s^{-1}}$.
\textbf{6.} $1670/(2\times 10^{-3}) = \qty{8.3e5}{L.mol^{-1}.s^{-1}}$: mil
veces por debajo del límite de difusión; solo una colisión de cada mil es
productiva.
\textbf{7.} $0.95 = [S]/(2 + [S])$: $[S] = 19\,K_m = \qty{38}{mmol/L}$.
\textbf{8.} $1/v_0$: 9,0, 5,0, 3,0, 1,8, 1,4, 1,2. Pendiente $(9.0 -
1.2)/1.95 = 4.0$; ordenada en el origen $1.0$.
\textbf{9.} $V_{\max}^{\text{ap}} = \qty{1.0}{\micro mol/min}$ (sin
cambios); $K_m^{\text{ap}} = 4.0\times 1.0 = \qty{4.0}{mmol/L}$.
\textbf{10.} $K_m$ duplicada, $V_{\max}$ sin cambios: competitivo.
\textbf{11.} $4 = 2(1 + 1/K_i)$: $K_i = \qty{1.0}{mmol/L}$.
\textbf{12.} Con el \omterm{def:b1:enzymes:inhibitors}{inhibidor}, $v = V_{\max}[S]/(K_m(1 + [I]/K_i) +
[S])$;
reducir la velocidad a la mitad exige $K_m(1 + [I]/K_i) + [S] = 2(K_m +
[S])$, es decir, $[I] = K_i(K_m + [S])/K_m$: a \qty{2}{mmol/L}, $[I] =
1\times 4/2 = \qty{2}{mmol/L}$; a \qty{20}{mmol/L}, $[I] = 22/2 =
\qty{11}{mmol/L}$. Cuanto más \omterm{def:b1:enzymes:enzyme}{sustrato}, más \omterm{def:b1:enzymes:inhibitors}{inhibidor} hace falta.
\textbf{13.} Una molécula parecida al \omterm{def:b1:enzymes:enzyme}{sustrato} (o a su \omterm{prop:b1:enzymes:whatitdoes}{estado de
transición}) --- un análogo de éster que no puede hidrolizarse --- que se
une en el \omterm{def:b1:enzymes:enzyme}{centro activo}.
\textbf{14.} $1/v_0$: 10,0, 6,0, 4,0, 2,8, 2,4, 2,2; pendiente $4.0$,
ordenada en el origen $2.0$: $V_{\max}^{\text{ap}} = \qty{0.5}{\micro mol/min}$,
$K_m^{\text{ap}} = 4.0\times 0.5 = \qty{2.0}{mmol/L}$.
\textbf{15.} $V_{\max}$ reducida a la mitad, $K_m$ sin cambios: no
competitivo; $2 = 1 + 1/K_i$: $K_i = \qty{1.0}{mmol/L}$.
\textbf{16.} J se une a un sitio distinto del \omterm{def:b1:enzymes:enzyme}{centro activo}, tanto a la
\omterm{def:b1:enzymes:enzyme}{enzima} libre como al $ES$, con la misma afinidad: el \omterm{def:b1:enzymes:enzyme}{sustrato} no compite con
él, y a cualquier $[S]$ la mitad de las moléculas de \omterm{def:b1:enzymes:enzyme}{enzima} quedan
inactivadas.
\textbf{17.} Con J la velocidad se duplica (la mitad del doble de \omterm{def:b1:enzymes:enzyme}{enzima}
sigue activa); con I también se duplica (la velocidad es proporcional a la
\omterm{def:b1:enzymes:enzyme}{enzima} a cualquier $[S]$): duplicar la \omterm{def:b1:enzymes:enzyme}{enzima} nunca distingue a los
\omterm{def:b1:enzymes:inhibitors}{inhibidores}.
\textbf{18.} Velocidad a \qty{25}{\celsius}: $V_{\max}[S]/(K_m + [S])$ con
$V_{\max} = k_{\text{cat}}[E]$: por litro, $1670\times 10^{-8} =
\qty{1.67e-5}{mol/L/s}$; $\times 0.5/2.5 = \qty{3.3e-6}{mol/L/s}$; a
\qty{37}{\celsius}, $\times 2^{1.2} = 2.3$: \qty{7.7e-6}{mol/L/s}; en
\qty{1e-12}{L}: \qty{7.7e-18}{mol/s} $= \num{4.6e6}$ moléculas por
segundo.
\textbf{19.} No hace falta subirla: la velocidad (\num{4.6e6}) supera la
demanda (\num{3e6}); si tuviera que subir, la \omterm{def:b1:cell-unit-of-life:cell}{célula} podría elevar el
\omterm{def:b1:enzymes:enzyme}{sustrato}, fosforilar la \omterm{def:b1:enzymes:enzyme}{enzima} para rebajar su $K_m$, o fabricar más
\omterm{def:b1:enzymes:enzyme}{enzima}.
\textbf{20.} $K_m^{\text{ap}} = 2(1 + 0.4/0.2) = \qty{6}{mmol/L}$: la
velocidad queda en $\times 0.5/6.5$ en vez de $0.5/2.5$: el \qty{38}{\%} de
antes, \num{1.8e6} por segundo; el producto estrangula su propia síntesis
por debajo de la demanda, de modo que se consume y su nivel cae, lo que
libera a la \omterm{def:b1:enzymes:enzyme}{enzima}: un suministro que se autorregula.
\textbf{21.} Antes: $0.5/2.5 = 0.20$ de $V_{\max}$; después: $0.5/0.9 =
0.56$: casi el triple.
\textbf{22.} Casi nula: a pH 5 la histidina del \omterm{def:b1:enzymes:enzyme}{centro activo} está protonada
y no puede actuar como base, y los residuos ácidos del centro quedan
neutralizados; el estado de ionización de la \omterm{def:b1:enzymes:enzyme}{enzima}, y quizá su plegamiento,
son los equivocados. Está hecha para el \omterm{def:b1:eukaryotic-cell:organelle}{citosol}, no para el \omterm{def:b1:eukaryotic-cell:endomembrane}{lisosoma}.
\textbf{23.} $K_m$ cambia poco (la unión se debe sobre todo al resto del
centro); $k_{\text{cat}}$ se desploma en órdenes de magnitud, ya que la
histidina era la base catalítica que activaba el agua o la serina.
\textbf{24.} Cerca de $K_m$ la velocidad responde a los cambios de \omterm{def:b1:enzymes:enzyme}{sustrato}
(pendiente cercana a $V_{\max}/2K_m$); muy por encima, la \omterm{def:b1:enzymes:enzyme}{enzima} está
saturada y es insensible, y el exceso de \omterm{def:b1:enzymes:enzyme}{sustrato} sería una carga osmótica y
química. Trabajar cerca de $K_m$ hace que la vía sea controlable por el
suministro.
\textbf{25.} $K_m = \qty{2.0}{mmol/L}$, $V_{\max} = \qty{1.0}{\micro mol/min}$
(con \qty{10}{nmol/L} de \omterm{def:b1:enzymes:enzyme}{enzima}), $k_{\text{cat}} = \qty{1670}{s^{-1}}$; I
es competitivo, $K_i = \qty{1.0}{mmol/L}$; J es no competitivo,
$K_i =
\qty{1.0}{mmol/L}$.
\end{solution}
